\chapter{TESTING AND VERIFICATION}

\section{Test Cases}
Testing was performed through both manual execution and automated regression
tests. Manual runs were used to inspect the full pipeline from source code to
output, while automated tests were added for the semantic and runtime behaviour
of the implemented type checker and interpreter.

\begin{table}[H]
\centering
\caption{Representative Test Cases}
\begin{tabular}{p{3.5cm}p{5cm}p{4cm}}
\toprule
Category & Program Behaviour & Expected Result \\
\midrule
Arithmetic & Evaluate mixed numeric expressions & Correct value and inferred type \\
If-then-else & Choose one branch from a Boolean condition & Only matching branch runs \\
While-loop & Print during repeated execution & One output line per iteration \\
Functions & Infer parameters and return type & Valid call and correct return value \\
Type mismatch & Reassign incompatible type & Type error before runtime \\
\bottomrule
\end{tabular}
\end{table}

\subsection{Arithmetic Expressions}
Arithmetic expressions were tested using both integer-only and mixed
integer-float programs. These tests verified precedence and numeric promotion.
Expressions such as \texttt{x + y * 2} confirmed that multiplication is applied
before addition.

\subsection{Boolean Expressions}
Boolean expressions were tested inside \texttt{if} and \texttt{while}
conditions. The checker verified valid comparison operands, and the interpreter
used the resulting Boolean values for control flow.

\subsection{Assignment Statements}
Assignments were tested for both first-use inference and reassignment checks.
For example, a variable assigned \texttt{5} and later assigned
\texttt{"hello"} correctly triggers a type error.

\subsection{If-Then-Else}
Conditional programs were used to verify that exactly one branch is executed and
that non-Boolean conditions are rejected before runtime.

\subsection{While-Loop}
While-loops were tested with a countdown example. This confirmed that the loop
body executes repeatedly and that \texttt{print()} emits output progressively on
each iteration rather than only once at the end.

\subsection{Functions}
Function tests covered parameter passing, local execution, and return values.
The implementation infers parameter types from the first call and then enforces
that signature for later calls.

\subsection{Print}
The built-in \texttt{print()} function was tested with all supported runtime
types. The output includes the printed value together with its type, such as
\texttt{29.5 : Float} and \texttt{"plc" : String}.

\section{Type Error Detection}
The checker was verified with invalid programs involving incompatible
assignment, invalid arithmetic operands, wrong function arguments, and
non-Boolean conditions. In each case the system stopped before execution and
reported a source-positioned type error.

\section{Automated Test Suite}
Automated regression tests were added in
\texttt{tests/test\_member3\_pipeline.py} using the \texttt{unittest}
framework. The suite currently covers:

\begin{itemize}
	\item progressive output during while-loop execution,
	\item assignment type mismatch detection,
	\item function parameter and return type inference.
\end{itemize}

The tests can be run using:

\begin{lstlisting}
python3 -m unittest discover -s tests -v
\end{lstlisting}

\section{Error Handling}
The language system reports errors at three main stages:

\begin{itemize}
	\item lexical errors for malformed characters or unterminated strings,
	\item parse errors for invalid syntax such as missing semicolons,
	\item type errors for programs that violate static typing rules.
\end{itemize}

Runtime errors are also reported with source positions whenever execution
reaches an invalid state such as an undefined variable reference.
